\documentclass[12pt,a4paper]{article}
\usepackage{../../../myconfig}

\begin{document}

% =========================================================================
% TRANG 1: NỀN TẢNG ĐƯỜNG THẲNG TRONG KHÔNG GIAN (KHÓA GỌN 3 CÂU NHẬN BIẾT)
% =========================================================================
\titlebox{Chủ đề: Oxyz}{Phương trình đường thẳng}{Oxyz: Phương trình đường thẳng}

\vspace{-0.25cm}

\begin{theorybox}
    \textbf{\textcolor{deepblue}{\faLightbulb\hskip 6pt NỀN TẢNG: CẦN GÌ ĐỂ VIẾT PHƯƠNG TRÌNH ĐƯỜNG THẲNG?}}
    \par\vspace{4pt}
    \begin{minipage}[c]{0.65\linewidth}
        Để xác định duy nhất đường thẳng $d$, ta luôn cần đủ \textbf{2 yếu tố}:
        \[
            \begin{cases} 
                \textbf{1 Điểm đi qua:} & M_0(x_0; y_0; z_0) \\ 
                \textbf{1 Véctơ chỉ phương:} & \vec{u} = (a; b; c) \neq \vec{0} \quad (\text{giá } /\!/ \text{ hoặc trùng } d)
            \end{cases}
        \]
        \textbf{1. Phương trình tham số của đường thẳng $d$ ($t \in \mathbb{R}$):}
        \[
            \begin{cases} x = x_0 + at \\ y = y_0 + bt \\ z = z_0 + ct \end{cases}
        \]
        \textbf{2. Phương trình chính tắc của đường thẳng $d$ (khi $abc \neq 0$):}
        \[
            \dfrac{x - x_0}{a} = \dfrac{y - y_0}{b} = \dfrac{z - z_0}{c}
        \]
    \end{minipage}%
    \hfill
    \begin{minipage}[c]{0.32\linewidth}
        \centering
        \begin{tikzpicture}[scale=0.92]
            % Đường thẳng d
            \draw[thick, draw=deepblue] (-0.3,-0.3) -- (3.5,2.1) node[right, font=\bfseries\small] {$d$};
            
            % Điểm M0
            \coordinate (M0) at (1.1, 0.54);
            \fill[deepblue] (M0) circle (2pt);
            \node[deepblue, below right=1pt, font=\small] at (M0) {$M_0(x_0; y_0; z_0)$};
            
            % Vector u
            \coordinate (Ustart) at (0.3, 1.4);
            \draw[->, >=stealth, line width=1.3pt, accentred] (Ustart) -- ++(1.8, 1.08) node[above left, font=\small] {$\vec{u}$};
            \draw[dashed, draw=accentred!50] ($(Ustart)+(-0.3,-0.18)$) -- ++(2.4, 1.44);
        \end{tikzpicture}
    \end{minipage}
\end{theorybox}

\vspace{0.15cm}


% Câu 1: Khai thác cả Điểm đi qua và VTCP từ dạng Tham số
\questionLinesManual{0}{hỏi}{%
Trong không gian $Oxyz$, cho đường thẳng $d\colon \begin{cases} x = 2 - 3t \\ y = 1 + 2t \\ z = -4 + t \end{cases}$. Đường thẳng $d$ đi qua điểm nào và có một véctơ chỉ phương là véctơ nào dưới đây?
\begin{tasks}(2)
    \task Đi qua $M(2; 1; -4)$ và $\vec{u} = (-3; 2; 1)$.
    \task Đi qua $N(-3; 2; 1)$ và $\vec{u} = (2; 1; -4)$.
    \task Đi qua $P(-2; -1; 4)$ và $\vec{u} = (-3; 2; 1)$.
    \task Đi qua $Q(2; 1; -4)$ và $\vec{u} = (3; 2; 1)$.
\end{tasks}
}{}

\vspace{0.25cm}

% Câu 2: Khai thác cả Điểm đi qua và VTCP từ dạng Chính tắc
\questionLinesManual{0}{hỏi}{%
Trong không gian $Oxyz$, cho đường thẳng $d\colon \dfrac{x - 1}{2} = \dfrac{y + 2}{-3} = \dfrac{z - 4}{1}$. Đường thẳng $d$ đi qua điểm nào và có một véctơ chỉ phương là véctơ nào dưới đây?
}{%
\begin{tasks}(2)
    \task Đi qua $P(1; -2; 4)$ và $\vec{u} = (2; -3; 1)$.
    \task Đi qua $M(-1; 2; -4)$ và $\vec{u} = (2; -3; 1)$.
    \task Đi qua $P(1; -2; 4)$ và $\vec{u} = (2; 3; 1)$.
    \task Đi qua $Q(-1; 2; -4)$ và $\vec{u} = (-2; 3; 1)$.
\end{tasks}
}

\newpage


\newpage

% =========================================================================
% TRANG 2: DẠNG 1 (QUA 1 ĐIỂM + VTCP)
% =========================================================================
\dangbox{Dạng 1}{Viết phương trình đường thẳng qua một điểm và có VTCP}

\begin{theorybox}
    \begin{minipage}[c]{0.58\linewidth}
        Đường thẳng $d$ xác định bởi:
        \[
            d\colon \begin{cases} 
                \text{Qua } M(x_0; y_0; z_0) \\ 
                \text{VTCP } \vec{u}_d = (a; b; c) 
            \end{cases}
        \]
        \begin{itemize}
            \item \textbf{Phương trình tham số ($t \in \mathbb{R}$):}
            \[
                \begin{cases} 
                    x = x_0 + at \\ 
                    y = y_0 + bt \\ 
                    z = z_0 + ct 
                \end{cases}
            \]
            \item \textbf{Phương trình chính tắc (khi $abc \neq 0$):}
            \[
                \dfrac{x - x_0}{a} = \dfrac{y - y_0}{b} = \dfrac{z - z_0}{c}
            \]
        \end{itemize}
    \end{minipage}%
    \hfill
    \begin{minipage}[c]{0.40\linewidth}
        \centering
        \begin{tikzpicture}[scale=1.0]
            % Đường thẳng d
            \draw[thick, draw=deepblue] (0,0.2) -- (4.2,2.2) node[right, font=\bfseries\small] {$d$};
            
            % Điểm M0
            \coordinate (M) at (1.5, 0.91);
            \fill[deepblue] (M) circle (2pt);
            \node[deepblue, below right=1pt, font=\small] at (M) {$M(x_0;y_0;z_0)$};
            
            % Vector u_d
            \draw[->, >=stealth, line width=1.3pt, accentred] (2.4, 1.34) -- ++(1.4, 0.67) node[above left, font=\small] {$\vec{u}_d$};
        \end{tikzpicture}
    \end{minipage}
\end{theorybox}

\vspace{0.25cm}

\questionLinesManual{7}{hỏi}{%
Trong không gian $Oxyz$, viết phương trình tham số và phương trình chính tắc của đường thẳng $d$ đi qua điểm $M(2; -1; 3)$ và có véctơ chỉ phương $\vec{u} = (3; 1; -2)$.
}{}

\vspace{0.35cm}

\questionLinesManual{7}{hỏi}{%
Trong không gian $Oxyz$, viết phương trình tham số của đường thẳng $\Delta$ đi qua gốc tọa độ $O(0;0;0)$ và nhận véctơ $\vec{u} = (2; -3; 5)$ làm véctơ chỉ phương.
}{}

\newpage

% =========================================================================
% TRANG 3: DẠNG 2 (ĐI QUA HAI ĐIỂM)
% =========================================================================
\dangbox{Dạng 2}{Viết phương trình đường thẳng đi qua hai điểm phân biệt $A, B$}

\begin{theorybox}
    \begin{minipage}[c]{0.58\linewidth}
        \[
            d\colon \begin{cases} 
                \text{Qua } A \text{ (hoặc } B\text{)} \\ 
                \text{VTCP } \vec{u}_d = \vec{AB} 
            \end{cases}
        \]
    \end{minipage}%
    \hfill
    \begin{minipage}[c]{0.40\linewidth}
        \centering
        \begin{tikzpicture}[scale=1.0]
            \draw[thick, draw=deepblue] (0,0.3) -- (4.2,2.0) node[right, font=\bfseries\small] {$d$};
            
            \coordinate (A) at (1.0, 0.7);
            \coordinate (B) at (3.0, 1.51);
            
            \fill[deepblue] (A) circle (2pt) node[above left=1pt, font=\small] {$A$};
            \fill[deepblue] (B) circle (2pt) node[below right=1pt, font=\small] {$B$};
            
            \draw[->, >=stealth, line width=1.3pt, accentred] (A) -- (B) node[midway, above left, font=\small] {$\vec{u}_d = \vec{AB}$};
        \end{tikzpicture}
    \end{minipage}
\end{theorybox}

\vspace{0.25cm}

\questionLinesManual{10}{hỏi}{%
Trong không gian $Oxyz$, viết phương trình tham số và chính tắc của đường thẳng $d$ đi qua hai điểm $A(1; 2; -1)$ và $B(2; -1; 1)$.
}{}

\vspace{0.35cm}

\questionLinesManual{12}{hỏi}{%
Trong không gian $Oxyz$, cho tam giác $ABC$ có $A(2; 1; 0)$, $B(0; -1; 2)$ và $C(4; 3; -2)$. Viết phương trình đường trung tuyến kẻ từ đỉnh $A$ của tam giác $ABC$.
}{}

\newpage

% =========================================================================
% TRANG 4: DẠNG 3 (SONG SONG VỚI ĐƯỜNG THẲNG ĐÃ BIẾT)
% =========================================================================
\dangbox{Dạng 3}{Viết phương trình đường thẳng qua $M$ và song song với $\Delta$}

\begin{theorybox}
    \begin{minipage}[c]{0.58\linewidth}
        \[
            d\colon \begin{cases} 
                \text{Qua } M \\ 
                \text{VTCP } \vec{u}_d = \vec{u}_\Delta 
            \end{cases}
        \]
    \end{minipage}%
    \hfill
    \begin{minipage}[c]{0.40\linewidth}
        \centering
        \begin{tikzpicture}[scale=1.0]
            % Đường thẳng Delta
            \draw[thick, draw=softgray!80] (0,0.4) -- (3.8,1.9) node[right, font=\bfseries\small] {$\Delta$};
            \draw[->, >=stealth, line width=1.1pt, deepblue] (1.2, 0.87) -- ++(1.4, 0.55) node[above left, font=\small] {$\vec{u}_\Delta$};
            
            % Đường thẳng d
            \draw[thick, draw=deepblue] (0,1.6) -- (3.8,3.1) node[right, font=\bfseries\small] {$d$};
            \coordinate (M) at (1.0, 1.99);
            \fill[deepblue] (M) circle (2pt) node[above left=1pt, font=\small] {$M$};
            \draw[->, >=stealth, line width=1.3pt, accentred] (M) -- ++(1.4, 0.55) node[above left, font=\small] {$\vec{u}_d = \vec{u}_\Delta$};
        \end{tikzpicture}
    \end{minipage}
\end{theorybox}

\vspace{0.25cm}

\questionLinesManual{10}{hỏi}{%
Trong không gian $Oxyz$, viết phương trình tham số của đường thẳng $d$ đi qua điểm $A(2; 1; -3)$ và song song với đường thẳng $\Delta\colon \dfrac{x - 1}{2} = \dfrac{y + 3}{-1} = \dfrac{z}{4}$.
}{}

\vspace{0.35cm}

\questionLinesManual{10}{hỏi}{%
Trong không gian $Oxyz$, viết phương trình tham số của đường thẳng $d$ đi qua điểm $M(1; -2; 4)$ và song song với trục hoành $Ox$.
}{}

\newpage

% =========================================================================
% TRANG 5: DẠNG 4 (VUÔNG GÓC VỚI MẶT PHẲNG)
% =========================================================================
\dangbox{Dạng 4}{Viết phương trình đường thẳng qua $M$ và vuông góc với mặt phẳng $(P)$}

\begin{theorybox}
    \begin{minipage}[c]{0.58\linewidth}
        \[
            d\colon \begin{cases} 
                \text{Qua } M \\ 
                \text{VTCP } \vec{u}_d = \vec{n}_{(P)} 
            \end{cases}
        \]
    \end{minipage}%
    \hfill
    \begin{minipage}[c]{0.40\linewidth}
        \centering
        \begin{tikzpicture}[scale=0.95]
            % Mặt phẳng (P)
            \draw[thick, fill=white, draw=deepblue] (0,0) -- (3.8,0) -- (4.8,1.2) -- (1.0,1.2) -- cycle;
            \draw[deepblue, thin] (0.4,0) arc (0:45:0.4);
            \node[deepblue, font=\bfseries\small] at (0.65, 0.22) {$(P)$};
            
            % Đường thẳng d vuông góc đâm xuyên
            \coordinate (H) at (2.4, 0.6);
            \draw[thick, deepblue] (H) -- (2.4, 2.3) node[above, font=\bfseries\small] {$d$};
            \draw[dashed, thick, deepblue] (H) -- (2.4, -0.3);
            
            \coordinate (M) at (2.4, 1.8);
            \fill[deepblue] (M) circle (1.8pt) node[right=2pt, font=\small] {$M$};
            
            \draw[->, >=stealth, line width=1.3pt, accentred] (H) -- ++(0, 1.0) node[left=1pt, font=\small] {$\vec{u}_d = \vec{n}_P$};
            \draw[accentred] (2.4, 0.8) -- (2.6, 0.8) -- (2.6, 0.6);
        \end{tikzpicture}
    \end{minipage}
\end{theorybox}

\vspace{0.25cm}

\questionLinesManual{8}{hỏi}{%
Trong không gian $Oxyz$, viết phương trình chính tắc của đường thẳng $d$ đi qua gốc tọa độ $O$ và vuông góc với mặt phẳng $(P)\colon 2x - 3y + z - 5 = 0$.
}{}

\vspace{0.35cm}

\questionLinesManual{12}{hỏi}{%
Trong không gian $Oxyz$, cho tam giác $ABC$ có ba đỉnh $A(1; 0; 0)$, $B(0; 2; 0)$, $C(0; 0; 3)$. Viết phương trình tham số của đường thẳng $d$ đi qua trực tâm $H$ của tam giác $ABC$ và vuông góc với mặt phẳng $(ABC)$.
}{}

\newpage

% =========================================================================
% TRANG 6: DẠNG 5 (GIAO TUYẾN CỦA HAI MẶT PHẲNG)
% =========================================================================
\dangbox{Dạng 5}{Viết phương trình đường thẳng là giao tuyến của hai mặt phẳng}

\begin{theorybox}
    \begin{minipage}[c]{0.62\linewidth}
        \[
            d\colon \begin{cases} 
                \text{VTCP } \vec{u}_d = [\vec{n}_{(P)}, \vec{n}_{(Q)}] \\[8pt]
                \text{Điểm } A \in d\colon \text{Cho } z=0 \text{ (hoặc } x=0, y=0\text{)}, \\
                \phantom{\text{Điểm } A \in d\colon} \text{giải hệ } \begin{cases} (P) \\ (Q) \end{cases} \text{tìm 2 ẩn còn lại.}
            \end{cases}
        \]
    \end{minipage}%
    \hfill
    \begin{minipage}[c]{0.36\linewidth}
        \centering
        \begin{tikzpicture}[scale=0.95, line join=round, line cap=round]
            % Tọa độ các đỉnh chuẩn hình học không gian
            \coordinate (TL1) at (0, 2.1);     % Đỉnh trên trái (mặt sau)
            \coordinate (BL1) at (0, 0.5);     % Đáy dưới trái (mặt sau)
            \coordinate (TR)  at (2.4, 1.7);   % Đỉnh trên của d
            \coordinate (BR)  at (2.4, -0.2);  % Chân đáy của d

            \coordinate (TL2) at (0.7, 1.25);  % Đỉnh trên trái (mặt trước)
            \coordinate (BL2) at (0.7, -0.65); % Đáy dưới trái (mặt trước)

            % Giao điểm giữa cạnh đáy mặt sau (BL1--BR) và cạnh đứng mặt trước (TL2--BL2)
            % Cạnh đứng có hoành độ x = 0.7 => y = 0.5 + (0.7/2.4)*(-0.2 - 0.5) ≈ 0.3
            \coordinate (X) at (0.7, 0.296);

            % --- 1. MẶT PHẲNG PHÍA SAU ---
            \draw[line width=1.1pt, draw=deepblue] (TL1) -- (TR);
            \draw[line width=1.1pt, draw=deepblue] (TL1) -- (BL1);
            \draw[line width=1.1pt, draw=deepblue] (BL1) -- (X);        % Đoạn đáy nhìn thấy ngoài
            \draw[dashed, line width=1.0pt, draw=deepblue] (X) -- (BR); % Đoạn đáy khuất sau mặt trước

            % --- 2. MẶT PHẲNG PHÍA TRƯỚC (đè lên trước) ---
            \draw[line width=1.1pt, draw=deepblue] (TL2) -- (TR);
            \draw[line width=1.1pt, draw=deepblue] (TL2) -- (BL2);
            \draw[line width=1.1pt, draw=deepblue] (BL2) -- (BR);

            % --- 3. ĐƯỜNG GIAO TUYẾN d VÀ ĐIỂM A ---
            \draw[line width=1.5pt, draw=deepblue] (TR) -- (BR);
            
            \coordinate (A) at ($(TR)!0.48!(BR)$);
            \fill[deepblue] (A) circle (2pt);
            \node[right=2pt, font=\bfseries\small, deepblue] at (A) {$A$};
            \node[left=3pt, font=\itshape\Large, deepblue] at ($(TR)!0.62!(BR)$) {$d$};
        \end{tikzpicture}
    \end{minipage}
\end{theorybox}

\vspace{0.25cm}

\questionLinesManual{9}{hỏi}{%
Trong không gian $Oxyz$, viết phương trình tham số của đường thẳng $d$ là giao tuyến của hai mặt phẳng $(P)\colon x + 2y - z - 3 = 0$ và $(Q)\colon 2x - y + z - 1 = 0$.
}{}

\vspace{0.35cm}

\questionLinesManual{9}{hỏi}{%
Trong không gian $Oxyz$, cho hai mặt phẳng $(\alpha)\colon 2x - y + 3z - 6 = 0$ và $(\beta)\colon x + 2y - z + 2 = 0$. Tìm một điểm thuộc giao tuyến của $(\alpha)$ và $(\beta)$ có cao độ $z = 0$, từ đó lập phương trình chính tắc của giao tuyến.
}{}

\newpage

% =========================================================================
% TRANG 7: DẠNG 6 (VUÔNG GÓC VỚI HAI ĐƯỜNG THẲNG)
% =========================================================================
\dangbox{Dạng 6}{Viết phương trình đường thẳng qua $M$ và vuông góc với $d_1, d_2$}

\begin{theorybox}
    \begin{minipage}[c]{0.58\linewidth}
        \[
            d\colon \begin{cases} 
                \text{Qua } M \\ 
                \text{VTCP } \vec{u}_d = [\vec{u}_{d_1}, \vec{u}_{d_2}] 
            \end{cases}
        \]
    \end{minipage}%
    \hfill
    \begin{minipage}[c]{0.40\linewidth}
        \centering
        \begin{tikzpicture}[scale=0.95, line join=round, line cap=round, >=stealth]
    % Đường d nằm ngang
    \draw[thick, draw=deepblue] (0, 0) -- (4.8, 0) node[below left=2pt, font=\bfseries\small] {$d$};

    % --- ĐƯỜNG d1 (thẳng đứng vuông góc với d) ---
    \coordinate (H1) at (1.1, 0);
    \draw[thick, draw=softgray!90] (1.1, -0.9) -- (1.1, 1.7) node[below right=2pt, font=\small] at (1.1, -0.1) {$d_1$};
    
    % Ký hiệu góc vuông tại d1
    \draw[accentred, thin] (0.9, 0) -- (0.9, 0.2) -- (1.1, 0.2);
    
    % Vectơ u_{d1}
    \draw[->, line width=1.1pt, accentred] (1.25, 0) -- (1.25, 1.1) node[above right=-1pt, font=\footnotesize] {$\vec{u}_{d_1}$};

    % --- ĐƯỜNG d2 (nghiêng vuông góc với d) ---
    \coordinate (H2) at (3.3, 0);
    \draw[thick, draw=softgray!90] (2.9, -0.9) -- (3.9, 1.7) node[below left=2pt, font=\small] at (3.3, -0.1) {$d_2$};
    
    % Ký hiệu góc vuông tại d2 (xoay theo hướng xiên)
    \draw[accentred, thin] (3.08, 0) -- (3.14, 0.16) -- (3.36, 0.16);
    
    % Vectơ u_{d2}
    \draw[->, line width=1.1pt, accentred] (3.5, 0) -- (3.9, 0.95) node[above right=-1pt, font=\footnotesize] {$\vec{u}_{d_2}$};
\end{tikzpicture}
    \end{minipage}
\end{theorybox}

\vspace{0.25cm}

\questionLinesManual{10}{hỏi}{%
Trong không gian $Oxyz$, viết phương trình chính tắc của đường thẳng $\Delta$ đi qua điểm $M(1; -1; 2)$ đồng thời vuông góc với cả hai đường thẳng $d_1\colon \dfrac{x}{1} = \dfrac{y - 1}{2} = \dfrac{z + 1}{-1}$ và $d_2\colon \dfrac{x - 1}{2} = \dfrac{y}{1} = \dfrac{z - 2}{1}$.
}{}

\vspace{0.35cm}

\questionLinesManual{10}{hỏi}{%
Trong không gian $Oxyz$, viết phương trình đường thẳng $d$ đi qua điểm $A(2; 0; -1)$ đồng thời song song với hai mặt phẳng $(P)\colon 2x - y + z - 3 = 0$ và $(Q)\colon x + 2y - z + 5 = 0$.
}{}

\newpage

% =========================================================================
% DẠNG 7: QUA ĐIỂM, VUÔNG GÓC VÀ CẮT DELTA
% =========================================================================
\dangbox{Dạng 7}{Viết phương trình đường thẳng qua điểm, vuông góc và cắt $\Delta$}

\begin{theorybox}
    \begin{minipage}[c]{0.58\linewidth}
        Gọi $H = d \cap \Delta \implies H(t) \in \Delta$.
        \[
            d \perp \Delta \iff \vec{AH} \cdot \vec{u}_\Delta = 0 \implies \text{tìm } t \implies \text{tọa độ } H.
        \]
        \[
            d\colon \begin{cases} 
                \text{Qua } A \\ 
                \text{VTCP } \vec{u}_d = \vec{AH} 
            \end{cases}
        \]
    \end{minipage}%
    \hfill
    \begin{minipage}[c]{0.40\linewidth}
        \centering
        \begin{tikzpicture}[scale=0.95, >=stealth]
            \draw[thick, draw=deepblue] (0, 0.4) -- (3.6, 0.9) node[right, font=\bfseries\small] {$\Delta$};
            
            \coordinate (H) at (1.8, 0.65);
            \coordinate (A) at (1.6, 2.1);
            
            \draw[thick, draw=accentred] (A) -- (H);
            \draw[thick, draw=accentred] (H) -- ($(H)!0.4!(2.0, -0.6)$) node[below, font=\bfseries\small] {$d$};
            
            \fill[deepblue] (A) circle (2pt) node[above=2pt, font=\small] {$A$};
            \fill[accentred] (H) circle (2pt);
            \node[below right=2pt, font=\small, accentred] at (H) {$H$};
            
            \draw[accentred, thin] ($(H)+(-0.03, 0.22)$) -- ($(H)+(0.19, 0.25)$) -- ($(H)+(0.22, 0.03)$);
        \end{tikzpicture}
    \end{minipage}
\end{theorybox}

\vspace{0.25cm}

\questionLinesManual{0}{hỏi}{%
Trong không gian $Oxyz$, cho điểm $A(1; 0; -2)$ và đường thẳng $\Delta\colon \dfrac{x - 1}{2} = \dfrac{y + 1}{-1} = \dfrac{z}{1}$. Viết phương trình tham số của đường thẳng $d$ đi qua $A$, vuông góc và cắt $\Delta$.
}{}
\drawdashlinesfull

\vspace{0.35cm}

\questionLinesManual{0}{hỏi}{%
Trong không gian $Oxyz$, cho điểm $M(2; -1; 3)$ và đường thẳng $\Delta\colon \begin{cases} x = 1 + 2t \\ y = -2 + t \\ z = 1 - t \end{cases}$. Tìm tọa độ hình chiếu vuông góc $H$ của $M$ lên $\Delta$, từ đó viết phương trình chính tắc của đường thẳng đi qua $M$ và cắt $\Delta$ tại $H$.
}{}
\drawdashlinesfull

\newpage
% =========================================================================
% DẠNG 8: QUA M, CẮT d1 VÀ VUÔNG GÓC d2
% =========================================================================
\dangbox{Dạng 8}{Viết phương trình đường thẳng qua $M$, cắt $d_1$ và vuông góc $d_2$}

\begin{theorybox}
    \begin{minipage}[c]{0.58\linewidth}
        Gọi $H = d \cap d_1 \implies H(t) \in d_1$.
        \[
            d \perp d_2 \iff \vec{MH} \cdot \vec{u}_{d_2} = 0 \implies \text{tìm } t \implies \text{tọa độ } H.
        \]
        \[
            d\colon \begin{cases} 
                \text{Qua } M \\ 
                \text{VTCP } \vec{u}_d = \vec{MH} 
            \end{cases}
        \]
    \end{minipage}%
    \hfill
    \begin{minipage}[c]{0.40\linewidth}
        \centering
        \begin{tikzpicture}[scale=0.92, line join=round, line cap=round]
            \draw[thick, draw=deepblue] (-0.2, 0.8) -- (4.2, 0.8) node[right, font=\bfseries\small] {$d$};
            
            \coordinate (H) at (1.5, 0.8);
            \coordinate (M) at (2.7, 0.8);
            \fill[deepblue] (H) circle (2pt) node[above left=1pt, font=\small] {$H$};
            \fill[deepblue] (M) circle (2pt) node[above=2pt, font=\small] {$M$};
            \draw[->, >=stealth, line width=1.2pt, accentred] (M) -- (H) node[midway, below=1pt, font=\footnotesize] {$\vec{u}_d$};

            \draw[thick, draw=softgray!90] (0.8, -0.4) -- (2.2, 2.0) node[above, font=\small] {$d_1$};
            \draw[thick, draw=softgray!90] (3.6, -0.4) -- (3.6, 2.0) node[above, font=\small] {$d_2$};
            \draw[->, >=stealth, line width=1.1pt, accentred] (3.6, 0.8) -- (3.6, 0.0) node[right=1pt, font=\footnotesize] {$\vec{u}_{d_2}$};
            \draw[accentred, thin] (3.6, 0.98) -- (3.42, 0.98) -- (3.42, 0.8);
        \end{tikzpicture}
    \end{minipage}
\end{theorybox}

\vspace{0.25cm}

\questionLinesManual{0}{hỏi}{%
Trong không gian $Oxyz$, cho điểm $M(1; 0; 2)$ và hai đường thẳng $d_1\colon \dfrac{x - 1}{2} = \dfrac{y}{1} = \dfrac{z + 1}{-1}$, $d_2\colon \dfrac{x}{1} = \dfrac{y - 2}{-2} = \dfrac{z + 1}{3}$. Viết phương trình tham số của đường thẳng $d$ đi qua $M$, cắt $d_1$ và vuông góc với $d_2$.
}{}
\drawdashlinesfull

\vspace{0.35cm}

\questionLinesManual{0}{hỏi}{%
Trong không gian $Oxyz$, viết phương trình chính tắc của đường thẳng $d$ đi qua điểm $A(2; -1; 1)$, vuông góc với trục $Oz$ và cắt đường thẳng $\Delta\colon \begin{cases} x = 1 + 3t \\ y = -2 + t \\ z = 2 - t \end{cases}$.
}{}
\drawdashlinesfull

\newpage
% =========================================================================
% DẠNG 9: NẰM TRONG MẶT PHẲNG (P) VÀ CẮT CẢ d1, d2
% =========================================================================
\dangbox{Dạng 9}{Viết phương trình đường thẳng nằm trong $(P)$ và cắt cả $d_1, d_2$}

\begin{theorybox}
    \begin{minipage}[c]{0.58\linewidth}
        Tìm tọa độ 2 giao điểm:
        \[
            A = d_1 \cap (P), \quad B = d_2 \cap (P)
        \]
        \[
            d\colon \begin{cases} 
                \text{Qua } A \text{ (hoặc } B\text{)} \\ 
                \text{VTCP } \vec{u}_d = \vec{AB} 
            \end{cases}
        \]
    \end{minipage}%
    \hfill
    \begin{minipage}[c]{0.40\linewidth}
        \centering
        \begin{tikzpicture}[scale=0.92, line join=round, line cap=round]
            \draw[thick, fill=white, draw=deepblue] (0,0) -- (4.2,0) -- (5.2,1.5) -- (1.0,1.5) -- cycle;
            \node[deepblue, font=\bfseries\small] at (0.6, 0.3) {$(P)$};

            \coordinate (A) at (1.8, 0.7);
            \coordinate (B) at (3.6, 0.8);
            
            \draw[line width=1.3pt, draw=accentred] (1.0, 0.65) -- (4.4, 0.85) node[right, font=\bfseries\small] {$d$};
            \fill[deepblue] (A) circle (2pt) node[below=2pt, font=\small] {$A$};
            \fill[deepblue] (B) circle (2pt) node[below=2pt, font=\small] {$B$};

            % d1 nghiêng dốc sang phải
            \draw[thick, draw=softgray!90] (1.3, 1.9) -- (A);
            \draw[dashed, thick, draw=softgray!90] (A) -- (2.1, -0.1) node[below, font=\small] {$d_1$};

            % d2 nghiêng dốc sang trái (lệch hẳn góc so với d1, không còn song song)
            \draw[thick, draw=softgray!90] (3.9, 1.9) -- (B);
            \draw[dashed, thick, draw=softgray!90] (B) -- (3.3, -0.2) node[below, font=\small] {$d_2$};
        \end{tikzpicture}
    \end{minipage}
\end{theorybox}

\vspace{0.25cm}

\questionLinesManual{0}{hỏi}{%
Trong không gian $Oxyz$, cho mặt phẳng $(P)\colon x + y + z - 4 = 0$ và hai đường thẳng $d_1\colon \dfrac{x - 1}{2} = \dfrac{y - 2}{1} = \dfrac{z}{-1}$, $d_2\colon \dfrac{x - 3}{1} = \dfrac{y + 1}{-1} = \dfrac{z - 2}{2}$. Viết phương trình chính tắc của đường thẳng $d$ nằm trong $(P)$ đồng thời cắt cả $d_1$ và $d_2$.
}{}
\drawdashlinesfull

\vspace{0.35cm}

\questionLinesManual{0}{hỏi}{%
Trong không gian $Oxyz$, cho mặt phẳng tọa độ $(Oxy)$ và hai đường thẳng $\Delta_1\colon \begin{cases} x = 1 + t \\ y = 2 - t \\ z = 1 + 2t \end{cases}$, $\Delta_2\colon \begin{cases} x = -2 + 3t' \\ y = 1 - t' \\ z = -2 + t' \end{cases}$. Viết phương trình tham số của đường thẳng $d$ nằm trong $(Oxy)$ và cắt cả $\Delta_1, \Delta_2$.
}{}
\drawdashlinesfull

\newpage
% =========================================================================
% DẠNG 10: QUA M VÀ CẮT CẢ HAI ĐƯỜNG THẲNG d1, d2
% =========================================================================
\dangbox{Dạng 10}{Viết phương trình đường thẳng qua $M$ và cắt hai đường thẳng $d_1, d_2$}

\begin{theorybox}
    \begin{minipage}[c]{0.58\linewidth}
        Gọi $A = d \cap d_1 \implies A(t)$ và $B = d \cap d_2 \implies B(t')$.
        \[
            M, A, B \text{ thẳng hàng } \iff \vec{MA} = k \vec{MB} \implies \text{tìm } t, t' \implies A, B.
        \]
        \[
            d\colon \begin{cases} 
                \text{Qua } M \\ 
                \text{VTCP } \vec{u}_d = \vec{MA} \text{ (hoặc } \vec{AB}\text{)} 
            \end{cases}
        \]
    \end{minipage}%
    \hfill
    \begin{minipage}[c]{0.40\linewidth}
        \centering
        \begin{tikzpicture}[scale=0.92, line join=round, line cap=round]
    % Đường thẳng d xuyên qua
    \draw[thick, draw=deepblue] (0, 0.4) -- (4.4, 1.8) node[right, font=\bfseries\small] {$d$};
    
    % Tọa độ các điểm trên d
    \coordinate (A) at (1.1, 0.75);
    \coordinate (M) at (2.2, 1.1);
    \coordinate (B) at (3.3, 1.45);
    
    \fill[deepblue] (M) circle (2pt) node[above=2pt, font=\small] {$M$};
    \fill[accentred] (A) circle (2pt) node[below right=1pt, font=\small] {$A$};
    \fill[accentred] (B) circle (2pt) node[below right=1pt, font=\small] {$B$};

    % d1 nghiêng sang trái (góc tù)
    \draw[thick, draw=softgray!90] (0.4, 2.0) -- (1.8, -0.5) node[below, font=\small] {$d_1$};
    
    % d2 nghiêng ngược lại sang phải (rõ ràng chéo nhau, không hề song song)
    \draw[thick, draw=softgray!90] (3.9, 2.2) -- (2.7, 0.7) node[below left=-1pt, font=\small] {$d_2$};
\end{tikzpicture}
    \end{minipage}
\end{theorybox}

\vspace{0.25cm}

\questionLinesManual{0}{hỏi}{%
Trong không gian $Oxyz$, cho điểm $M(1; 1; 1)$ và hai đường thẳng $d_1\colon \dfrac{x - 2}{1} = \dfrac{y + 1}{2} = \dfrac{z}{-1}$, $d_2\colon \dfrac{x + 1}{2} = \dfrac{y - 1}{1} = \dfrac{z - 1}{1}$. Viết phương trình tham số của đường thẳng $d$ đi qua điểm $M$ đồng thời cắt cả hai đường thẳng $d_1$ và $d_2$.
}{}
\drawdashlinesfull

\vspace{0.35cm}

\questionLinesManual{0}{hỏi}{%
Trong không gian $Oxyz$, viết phương trình đường thẳng $d$ đi qua gốc tọa độ $O$ đồng thời cắt cả hai đường thẳng $d_1\colon \begin{cases} x = 1 + t \\ y = 2 - t \\ z = 3t \end{cases}$ và $d_2\colon \begin{cases} x = 2 - 2t' \\ y = -1 + t' \\ z = 1 + 3t' \end{cases}$.
}{}
\drawdashlinesfull

\newpage
% =========================================================================
% DẠNG 11: SONG SONG VỚI DELTA VÀ CẮT d1, d2
% =========================================================================
\dangbox{Dạng 11}{Viết phương trình đường thẳng song song $\Delta$ và cắt $d_1, d_2$}

\begin{theorybox}
    \begin{minipage}[c]{0.58\linewidth}
        Gọi $A = d \cap d_1 \implies A(t)$ và $B = d \cap d_2 \implies B(t')$.
        \[
            d \mathbin{/\!/} \Delta \iff \vec{AB} = k \vec{u}_\Delta \implies \text{tìm } t, t' \implies \text{tọa độ } A.
        \]
        \[
            d\colon \begin{cases} 
                \text{Qua } A \\ 
                \text{VTCP } \vec{u}_d = \vec{u}_\Delta 
            \end{cases}
        \]
    \end{minipage}%
    \hfill
    \begin{minipage}[c]{0.40\linewidth}
        \centering
        \begin{tikzpicture}[scale=0.92, line join=round, line cap=round]
            \draw[thick, draw=softgray!90] (0.2, 2.2) -- (3.8, 2.2) node[right, font=\bfseries\small] {$\Delta$};
            \draw[->, >=stealth, line width=1.1pt, accentred] (1.3, 2.2) -- ++(1.2, 0) node[above=1pt, font=\footnotesize] {$\vec{u}_\Delta$};

            \draw[thick, draw=deepblue] (0.2, 1.0) -- (4.2, 1.0) node[right, font=\bfseries\small] {$d$};
            \coordinate (A) at (1.2, 1.0);
            \coordinate (B) at (3.1, 1.0);
            \fill[accentred] (A) circle (2pt) node[below left=1pt, font=\small] {$A$};
            \fill[accentred] (B) circle (2pt) node[below right=1pt, font=\small] {$B$};

            \draw[thick, draw=softgray!90] (0.5, 1.8) -- (1.8, 0.1) node[below, font=\small] {$d_1$};
            \draw[thick, draw=softgray!90] (3.6, 1.9) -- (2.6, 0.1) node[below, font=\small] {$d_2$};
        \end{tikzpicture}
    \end{minipage}
\end{theorybox}

\vspace{0.25cm}

\questionLinesManual{0}{hỏi}{%
Trong không gian $Oxyz$, cho đường thẳng $\Delta\colon \dfrac{x - 1}{1} = \dfrac{y + 2}{-1} = \dfrac{z}{2}$ và hai đường thẳng $d_1\colon \dfrac{x - 2}{2} = \dfrac{y - 1}{1} = \dfrac{z + 1}{-1}$, $d_2\colon \dfrac{x}{1} = \dfrac{y + 1}{2} = \dfrac{z - 2}{1}$. Viết phương trình tham số của đường thẳng $d$ song song với $\Delta$ đồng thời cắt cả $d_1$ và $d_2$.
}{}
\drawdashlinesfull

\vspace{0.35cm}

\questionLinesManual{0}{hỏi}{%
Trong không gian $Oxyz$, viết phương trình chính tắc của đường thẳng $d$ song song với trục $Oz$ đồng thời cắt cả hai đường thẳng $d_1\colon \begin{cases} x = 2 + t \\ y = 1 - t \\ z = 3t \end{cases}$ và $d_2\colon \begin{cases} x = 1 - 2t' \\ y = 3 + t' \\ z = 2 - t' \end{cases}$.
}{}
\drawdashlinesfull

\newpage
% =========================================================================
% DẠNG 12: ĐƯỜNG VUÔNG GÓC CHUNG CỦA HAI ĐƯỜNG THẲNG CHÉO NHAU
% =========================================================================
\dangbox{Dạng 12}{Lập phương trình đường vuông góc chung của hai đường thẳng chéo nhau}

\begin{theorybox}
    \begin{minipage}[c]{0.58\linewidth}
        Gọi $A = \Delta \cap d_1 \implies A(t)$ và $B = \Delta \cap d_2 \implies B(t')$.
        \[
            \begin{cases} 
                \vec{AB} \cdot \vec{u}_{d_1} = 0 \\ 
                \vec{AB} \cdot \vec{u}_{d_2} = 0 
            \end{cases} 
            \implies \text{tìm } t, t' \implies \text{tọa độ } A, B.
        \]
        \[
            \Delta\colon \begin{cases} 
                \text{Qua } A \\ 
                \text{VTCP } \vec{u}_\Delta = \vec{AB} 
            \end{cases}
        \]
    \end{minipage}%
    \hfill
    \begin{minipage}[c]{0.40\linewidth}
        \centering
        \begin{tikzpicture}[scale=0.92]
            \draw[thick, draw=deepblue] (0.2, 1.8) -- (3.8, 2.3) node[right, font=\bfseries\small] {$d_1$};
            \draw[thick, draw=deepblue] (0.2, 0.4) -- (3.8, -0.1) node[right, font=\bfseries\small] {$d_2$};
            
            \coordinate (A) at (2.0, 2.05);
            \coordinate (B) at (2.0, 0.15);
            
            \fill[deepblue] (A) circle (2pt) node[above=2pt, font=\small] {$A$};
            \fill[deepblue] (B) circle (2pt) node[below=2pt, font=\small] {$B$};
            
            \draw[line width=1.3pt, accentred] (A) -- (B) node[midway, right=2pt, font=\small] {$\Delta$};
            
            \draw[accentred] (2.0, 1.83) -- (2.22, 1.86) -- (2.22, 2.08);
            \draw[accentred] (2.0, 0.37) -- (2.22, 0.34) -- (2.22, 0.12);
        \end{tikzpicture}
    \end{minipage}
\end{theorybox}

\vspace{0.25cm}

\questionLinesManual{15}{hỏi}{%
Trong không gian $Oxyz$, cho hai đường thẳng chéo nhau $d_1\colon \dfrac{x - 1}{2} = \dfrac{y - 2}{1} = \dfrac{z + 1}{2}$ và $d_2\colon \dfrac{x + 1}{1} = \dfrac{y - 1}{1} = \dfrac{z - 2}{-1}$. Tìm tọa độ hai điểm $A \in d_1, B \in d_2$ sao cho $AB$ là đoạn vuông góc chung của $d_1$ và $d_2$.
}{}

\vspace{0.35cm}

\questionLinesManual{3}{hỏi}{%
Từ kết quả tìm được ở câu trên, hãy viết phương trình tham số của đường thẳng $\Delta$ là đường vuông góc chung của hai đường thẳng $d_1$ và $d_2$.
}{}

\newpage
% =========================================================================
% DẠNG 13: HÌNH CHIẾU VUÔNG GÓC CỦA ĐIỂM LÊN MẶT PHẲNG
% =========================================================================
\dangbox{Dạng 13}{Tìm hình chiếu vuông góc của điểm lên mặt phẳng}

\begin{theorybox}
    \begin{minipage}[c]{0.58\linewidth}
        Tìm hình chiếu vuông góc $H$ của điểm $M$ lên $(P)$:
        \begin{enumerate}
            \item Viết phương trình đường thẳng $\Delta$ qua $M$ và $\Delta \perp (P)$:
            \[
                \Delta\colon \begin{cases} \text{Qua } M \\ \text{VTCP } \vec{u}_\Delta = \vec{n}_{(P)} \end{cases}
            \]
            \item Tọa độ $H$ là giao điểm của $\Delta$ và $(P)$ ($H = \Delta \cap (P)$).
        \end{enumerate}
    \end{minipage}%
    \hfill
    \begin{minipage}[c]{0.40\linewidth}
        \centering
        \begin{tikzpicture}[scale=0.92, line join=round, line cap=round]
            % Mặt phẳng (P)
            \draw[thick, fill=white, draw=deepblue] (0,0) -- (3.8,0) -- (4.8,1.3) -- (1.0,1.3) -- cycle;
            \node[deepblue, font=\bfseries\small] at (0.6, 0.25) {$(P)$};

            % Tọa độ
            \coordinate (H) at (2.4, 0.65);
            \coordinate (M) at (2.4, 2.3);

            % Đường thẳng delta vuông góc đâm xuyên
            \draw[thick, draw=accentred] (M) -- (H);
            \draw[dashed, thick, draw=accentred] (H) -- (2.4, -0.3) node[below, font=\bfseries\small] {$\Delta$};

            % Điểm M và H
            \fill[deepblue] (M) circle (2pt) node[above right=1pt, font=\small] {$M$};
            \fill[accentred] (H) circle (2pt) node[below right=1pt, font=\small] {$H$};

            % Ký hiệu vuông góc
            \draw[accentred, thin] (2.4, 0.85) -- (2.6, 0.85) -- (2.6, 0.65);
        \end{tikzpicture}
    \end{minipage}
\end{theorybox}

\vspace{0.25cm}

\questionLinesManual{0}{hỏi}{%
Trong không gian $Oxyz$, cho điểm $M(1; -2; 3)$ và mặt phẳng $(P)\colon 2x - y + 2z + 1 = 0$. Tìm tọa độ hình chiếu vuông góc $H$ của điểm $M$ lên mặt phẳng $(P)$, từ đó tìm tọa độ điểm $M'$ đối xứng với $M$ qua mặt phẳng $(P)$.
}{}
\drawdashlinesfull

\vspace{0.35cm}

\questionLinesManual{0}{hỏi}{%
Trong không gian $Oxyz$, cho mặt phẳng $(\alpha)\colon x + 2y - z - 6 = 0$. Tìm tọa độ hình chiếu vuông góc của gốc tọa độ $O$ lên mặt phẳng $(\alpha)$.
}{}
\drawdashlinesfull

\newpage
% =========================================================================
% DẠNG 14: HÌNH CHIẾU VUÔNG GÓC CỦA ĐIỂM LÊN ĐƯỜNG THẲNG
% =========================================================================
\dangbox{Dạng 14}{Tìm hình chiếu vuông góc của điểm lên đường thẳng}

\begin{theorybox}
    \begin{minipage}[c]{0.58\linewidth}
        Tìm hình chiếu vuông góc $H$ của điểm $M$ lên đường thẳng $d$:
        \[
            \begin{cases}
                H \in d \implies H(t) \\
                MH \perp d \iff \vec{MH} \cdot \vec{u}_d = 0
            \end{cases}
        \]
        Giải tìm $t \implies$ Tọa độ hình chiếu $H$.
    \end{minipage}%
    \hfill
    \begin{minipage}[c]{0.40\linewidth}
        \centering
        \begin{tikzpicture}[scale=0.92, line join=round, line cap=round]
            % Đường thẳng d nằm nghiêng nhẹ
            \draw[thick, draw=deepblue] (0.2, 0.4) -- (4.0, 1.0) node[right, font=\bfseries\small] {$d$};

            \coordinate (H) at (2.0, 0.7);
            \coordinate (M) at (1.8, 2.2);

            % Đoạn MH vuông góc
            \draw[thick, draw=accentred] (M) -- (H);

            % Điểm M và H
            \fill[deepblue] (M) circle (2pt) node[above=2pt, font=\small] {$M$};
            \fill[accentred] (H) circle (2pt);
            \node[below right=2pt, font=\small, accentred] at (H) {$H$};

            % Ký hiệu vuông góc xoay theo phương đường d
            \draw[accentred, thin] ($(H)+(-0.03, 0.22)$) -- ($(H)+(0.19, 0.25)$) -- ($(H)+(0.22, 0.03)$);
        \end{tikzpicture}
    \end{minipage}
\end{theorybox}

\vspace{0.25cm}

\questionLinesManual{0}{hỏi}{%
Trong không gian $Oxyz$, cho điểm $A(2; 1; -1)$ và đường thẳng $d\colon \dfrac{x - 1}{2} = \dfrac{y + 2}{1} = \dfrac{z - 3}{-1}$. Tìm tọa độ hình chiếu vuông góc $H$ của điểm $A$ lên đường thẳng $d$.
}{}
\drawdashlinesfull

\vspace{0.35cm}

\questionLinesManual{0}{hỏi}{%
Trong không gian $Oxyz$, cho điểm $M(1; 3; -2)$ và đường thẳng $\Delta\colon \begin{cases} x = 2 + t \\ y = -1 - t \\ z = 1 + 2t \end{cases}$. Tìm tọa độ điểm $M'$ đối xứng với điểm $M$ qua đường thẳng $\Delta$.
}{}
\drawdashlinesfull

\newpage
% =========================================================================
% CHUYÊN ĐỀ 1: CÔNG THỨC KHOẢNG CÁCH
% =========================================================================
\dangbox{Chuyên đề}{Công thức tính khoảng cách liên quan đến đường thẳng}

\begin{theorybox}
    \textbf{1. Khoảng cách từ điểm $A$ đến đường thẳng $d$:}
    \[
        d\colon \begin{cases} \text{Qua } M \\ \text{VTCP } \vec{u}_d \end{cases} 
        \implies d(A, d) = \frac{\left| [\vec{AM}, \vec{u}_d] \right|}{|\vec{u}_d|}
    \]
    
    \vspace{0.15cm}
    \textbf{2. Khoảng cách giữa hai đường thẳng $d_1$ và $d_2$:}
    \[
        d_1\colon \begin{cases} \text{Qua } M \\ \text{VTCP } \vec{u}_1 \end{cases}, \quad 
        d_2\colon \begin{cases} \text{Qua } N \\ \text{VTCP } \vec{u}_2 \end{cases}
    \]
    \begin{itemize}
        \item \textbf{Trường hợp 1 ($d_1 \mathbin{/\!/} d_2$):} $d(d_1, d_2) = d(M, d_2) = \dfrac{\left| [\vec{MN}, \vec{u}_2] \right|}{|\vec{u}_2|}$.
        \item \textbf{Trường hợp 2 ($d_1, d_2$ chéo nhau):} Đặt $\vec{u} = [\vec{u}_1, \vec{u}_2] \implies d(d_1, d_2) = \dfrac{\left| \vec{MN} \cdot \vec{u} \right|}{|\vec{u}|}$.
    \end{itemize}
\end{theorybox}

\vspace{0.25cm}

% Câu 1: Điểm đến đường thẳng
\questionLinesManual{14}{hỏi}{%
Trong không gian $Oxyz$, cho điểm $A(2; 1; -1)$ và đường thẳng $d\colon \dfrac{x - 1}{2} = \dfrac{y + 1}{1} = \dfrac{z}{-2}$. Tính khoảng cách từ điểm $A$ đến đường thẳng $d$.
}{}

\vspace{0.35cm}

% Câu 2: Hai đường thẳng song song
\questionLinesManual{10}{hỏi}{%
Trong không gian $Oxyz$, cho hai đường thẳng:
\[
    d_1\colon \frac{x - 1}{2} = \frac{y + 2}{-1} = \frac{z - 3}{1} \quad \text{và} \quad 
    d_2\colon \frac{x + 2}{2} = \frac{y - 1}{-1} = \frac{z + 1}{1}
\]
Tính khoảng cách giữa hai đường thẳng $d_1$ và $d_2$.
}{}

\vspace{0.35cm}

% Câu 3: Hai đường thẳng chéo nhau
\questionLinesManual{10}{hỏi}{%
Trong không gian $Oxyz$, cho hai đường thẳng:
\[
    d_1\colon \begin{cases} x = 1 + t \\ y = 2 - t \\ z = -1 + 2t \end{cases} \quad \text{và} \quad 
    d_2\colon \frac{x - 2}{2} = \frac{y + 1}{1} = \frac{z - 1}{-1}
\]
Chứng minh $d_1$ và $d_2$ chéo nhau, từ đó tính khoảng cách giữa $d_1$ và $d_2$.
}{}

\newpage
% =========================================================================
% CHUYÊN ĐỀ 2: VỊ TRÍ TƯƠNG ĐỐI GIỮA HAI ĐƯỜNG THẲNG
% =========================================================================
\dangbox{Chuyên đề}{Vị trí tương đối giữa hai đường thẳng}

\begin{theorybox}
    Cho hai đường thẳng:
    \[
        d_1\colon \begin{cases} \text{Qua } M \\ \text{VTCP } \vec{u}_1 \end{cases} \quad \text{và} \quad 
        d_2\colon \begin{cases} \text{Qua } N \\ \text{VTCP } \vec{u}_2 \end{cases}.
        \quad \text{Đặt } \vec{u} = [\vec{u}_1, \vec{u}_2]
    \]
    
    \vspace{0.1cm}
    \centering
    \begin{tikzpicture}[>=stealth, font=\small]
        % Gốc: Tích có hướng hai VTCP
        \node[draw=deepblue, thick, rounded corners=4pt, fill=white, inner sep=5pt] (root) at (0, 0) {$[\vec{u}_1, \vec{u}_2]$};

        % Nhánh trên: = 0 (cùng phương)
        \node[draw=deepblue, rounded corners=4pt, fill=white, inner sep=4pt] (node_up) at (4.5, 1.3) {$[\vec{MN}, \vec{u}_1]$};
        \draw[thick, draw=deepblue] (root) -- (1.5, 1.3) node[midway, above left=1pt, font=\footnotesize\bfseries, accentred] {$= \vec{0}$} -- (node_up);

        % Nhánh con nhánh trên: Trùng / Song song
        \coordinate (end_trung) at (8.0, 1.9);
        \coordinate (end_ss) at (8.0, 0.7);
        \draw[thick, draw=deepblue] (node_up) -- (6.3, 1.9) node[midway, above=1pt, font=\footnotesize\bfseries, accentred] {$= \vec{0}$} -- (end_trung) 
            node[right, font=\bfseries] {$d_1 \equiv d_2$};
        \draw[thick, draw=deepblue] (node_up) -- (6.3, 0.7) node[midway, below=1pt, font=\footnotesize\bfseries, accentred] {$\neq \vec{0}$} -- (end_ss) 
            node[right, font=\bfseries] {$d_1 \mathbin{/\!/} d_2$};

        % Nhánh dưới: \neq 0 (không cùng phương)
        \node[draw=deepblue, rounded corners=4pt, fill=white, inner sep=4pt] (node_down) at (4.5, -1.3) {$\vec{MN} \cdot \vec{u}$};
        \draw[thick, draw=deepblue] (root) -- (1.5, -1.3) node[midway, below left=1pt, font=\footnotesize\bfseries, accentred] {$\neq \vec{0}$} -- (node_down);

        % Nhánh con nhánh dưới: Cắt nhau / Chéo nhau
        \coordinate (end_cat) at (8.0, -0.7);
        \coordinate (end_cheo) at (8.0, -1.9);
        \draw[thick, draw=deepblue] (node_down) -- (6.3, -0.7) node[midway, above=1pt, font=\footnotesize\bfseries, accentred] {$= 0$} -- (end_cat) 
            node[right, font=\bfseries] {$d_1 \text{ cắt } d_2$};
        \draw[thick, draw=deepblue] (node_down) -- (6.3, -1.9) node[midway, below=1pt, font=\footnotesize\bfseries, accentred] {$\neq 0$} -- (end_cheo) 
            node[right, font=\bfseries] {$d_1 \text{ chéo } d_2$};
    \end{tikzpicture}
\end{theorybox}

\vspace{0.2cm}

% Câu 1: Trùng nhau
\questionLinesManual{11}{hỏi}{%
Trong không gian $Oxyz$, xét vị trí tương đối giữa hai đường thẳng:
\[
    d_1\colon \frac{x - 1}{2} = \frac{y + 1}{-1} = \frac{z - 2}{3} \quad \text{và} \quad 
    d_2\colon \begin{cases} x = 3 - 4t \\ y = -2 + 2t \\ z = 5 - 6t \end{cases}
\]
}{}

\vspace{0.25cm}

% Câu 2: Song song
\questionLinesManual{11}{hỏi}{%
Trong không gian $Oxyz$, xét vị trí tương đối giữa hai đường thẳng:
\[
    d_1\colon \frac{x - 2}{1} = \frac{y}{2} = \frac{z + 1}{-1} \quad \text{và} \quad 
    d_2\colon \frac{x + 1}{-2} = \frac{y - 3}{-4} = \frac{z}{2}
\]
}{}

\vspace{0.25cm}

% Câu 3: Cắt nhau
\questionLinesManual{11}{hỏi}{%
Trong không gian $Oxyz$, xét vị trí tương đối và tìm tọa độ giao điểm (nếu có) của hai đường thẳng:
\[
    d_1\colon \begin{cases} x = 1 + 2t \\ y = -1 + t \\ z = 2 - t \end{cases} \quad \text{và} \quad 
    d_2\colon \frac{x - 3}{1} = \frac{y}{2} = \frac{z - 1}{1}
\]
}{}

\vspace{0.25cm}

% Câu 4: Chéo nhau
\questionLinesManual{11}{hỏi}{%
Trong không gian $Oxyz$, xét vị trí tương đối giữa hai đường thẳng:
\[
    d_1\colon \frac{x - 1}{1} = \frac{y + 2}{2} = \frac{z}{3} \quad \text{và} \quad 
    d_2\colon \frac{x + 2}{2} = \frac{y - 1}{1} = \frac{z - 1}{1}
\]
}{}

\end{document}